\documentclass[11pt,a4paper]{amsart}

\usepackage[margin=1in]{geometry}
\usepackage{amsmath,amssymb,amsthm}
\usepackage[hidelinks]{hyperref}
\emergencystretch=1em

\newtheorem{theorem}{Theorem}[section]
\newtheorem{proposition}[theorem]{Proposition}
\newtheorem{lemma}[theorem]{Lemma}
\newtheorem{corollary}[theorem]{Corollary}
\theoremstyle{definition}
\newtheorem{definition}[theorem]{Definition}
\theoremstyle{remark}
\newtheorem{remark}[theorem]{Remark}

\newcommand{\T}{\mathbb T}
\newcommand{\N}{\mathbb N}
\newcommand{\Z}{\mathbb Z}
\newcommand{\R}{\mathbb R}
\newcommand{\E}{\mathbb E}
\newcommand{\Pp}{\mathbb P}
\newcommand{\ind}{\mathbf 1}
\newcommand{\dist}{\operatorname{dist}}
\newcommand{\Var}{\operatorname{Var}}
\newcommand{\Cov}{\operatorname{Cov}}
\newcommand{\Leb}{\operatorname{Leb}}
\newcommand{\e}{\mathrm e}

\title{A Cauchy Limit for a Sawtooth Sum}
\date{}
\subjclass[2020]{11K38, 11A55, 37A50, 60F05}
\keywords{continued fractions, discrepancy sums, Cauchy law, stable limits,
Gauss map}

\begin{document}
\begin{abstract}
For $S_n(\alpha)=\sum_{k\le n}(1/2-\{k\alpha\})$, with $\alpha$ uniformly distributed on $(0,1)$, we prove that $S_n(\alpha)/\log n$ converges in distribution to the centered Cauchy law of scale $1/(2\pi)$.  An exact Euclidean-algorithm reciprocity formula reduces the sum to an alternating continued-fraction cost.  Its principal part is a Gauss digit marked by a rapidly oscillating torus coordinate, while the remaining carry term is bounded.  Local cylinder oscillation, a resonance decomposition, and a pullback carry construction give the required Poisson limit for large jumps and a weak law for the bounded remainder.  Consequently the limiting distribution function is $\frac12+\pi^{-1}\arctan(2\pi c)$.
\end{abstract}

\maketitle

\begin{center}
\textbf{AI-use disclosure.} The proposed answer was first obtained in a
one-shot ChatGPT response using OpenAI GPT-5.6-sol Pro.  The proof was
subsequently developed and gaps were addressed through iterative prompting of
that model.
OpenAI Codex (GPT-5.6-sol) was used to assemble, edit, compile, and perform
automated consistency audits of this manuscript.
\end{center}

\section{Introduction and statement of the result}

Write $\{x\}=x-\lfloor x\rfloor$ and let logarithms be natural.  For
\[
 S_n(\alpha)=\sum_{k=1}^{n}\left(\frac12-\{k\alpha\}\right),
 \qquad 0<\alpha<1,
\]
the question is whether $S_n(\alpha)/\log n$, regarded as a random variable on
the Lebesgue probability space $(0,1)$, has a limiting law.  This is
Erd\H{o}s Problem~\#1002, originating in \cite{Er64b}; see also
\cite{Bloom1002}.  Kesten's two-parameter result \cite{Ke60} provides the
closely related averaged-origin case.  The answer to the fixed-origin problem
is as follows.

\begin{theorem}\label{thm:main}
As $n\to\infty$,
\[
 \frac{S_n(\alpha)}{\log n}
 \ \Longrightarrow\
 \operatorname{Cauchy}\!\left(0,\frac1{2\pi}\right).
\]
Equivalently, for every $c\in\R$,
\[
 \lim_{n\to\infty}\Leb\left\{\alpha\in(0,1):
 \frac{S_n(\alpha)}{\log n}\le c\right\}
 =\frac12+\frac1\pi\arctan(2\pi c).
\]
\end{theorem}

Here $\operatorname{Cauchy}(0,\gamma)$ denotes the centered Cauchy law with
characteristic function $t\mapsto\exp(-\gamma|t|)$.

The proof has four parts.  First an exact reciprocity identity expresses
$S_n$ as an alternating sum along the Gauss map.  Second, continued-fraction
coordinates separate each summand into a heavy-tailed marked digit and a
bounded carry remainder.  Third, the marked digits satisfy a Poisson limit;
the main point is that the initial measure lies on the oscillating graph
$\alpha\mapsto(\alpha,\{n\alpha\})$, so ordinary Gauss mixing is insufficient.
We instead integrate Fourier modes on complete continued-fraction cylinders,
before and after their unique resonance scale.  Finally, the carry is realized
as a measurable graph over the natural extension of a Gauss--torus cocycle;
a reset argument and a two-block correlation estimate give a weak law for the
bounded remainder.

Throughout, constants denoted by $C,c$ may change from line to line.  A bound
$O_A(L^{-A})$ is uniform in the time indices satisfying the displayed
hypotheses.  We put $L=\log n$ whenever $n$ is fixed.

\section{Exact renormalization and continued-fraction coordinates}

Let $T:(0,1)\to(0,1)$ be the Gauss map,
\[
 Tx=\left\{\frac1x\right\},
 \qquad a_1(x)=\left\lfloor\frac1x\right\rfloor.
\]
Rational points will always be ignored, since they form a null set.

\begin{proposition}[Reciprocity]\label{prop:reciprocity}
For irrational $x\in(0,1)$ and $N\ge1$, set
\[
 m=\lfloor Nx\rfloor,\qquad u=\{Nx\}.
\]
Then
\begin{equation}\label{eq:reciprocity}
 S_N(x)=-S_m(Tx)+\Phi(x,u),
 \qquad
 \Phi(x,u)=\frac{u(1-u)}{2x}-\frac u2.
\end{equation}
\end{proposition}

\begin{proof}
Since $x$ is irrational,
\[
 S_N(x)=\frac N2-\frac{xN(N+1)}2+\sum_{k=1}^{N}\lfloor kx\rfloor.
\]
Counting lattice points below $j=kx$ in the other order gives
\[
 \sum_{k=1}^{N}\lfloor kx\rfloor
 =Nm-\sum_{j=1}^{m}\left\lfloor\frac jx\right\rfloor.
\]
If $a=\lfloor1/x\rfloor$, then $1/x=a+Tx$, and hence
\[
 \sum_{j=1}^{m}\left\lfloor\frac jx\right\rfloor
 =\frac{am(m+1)}2+\sum_{j=1}^{m}\lfloor jTx\rfloor.
\]
Substitution and the identity $m=Nx-u$ reduce the remaining polynomial
expression to $u(1-u)/(2x)-u/2$.
\end{proof}



Starting with $N_0=n$, $x_0=\alpha$, define
\begin{equation}\label{eq:descent}
 N_{j+1}=\lfloor N_jx_j\rfloor,
 \qquad x_{j+1}=Tx_j,
 \qquad u_j=\{N_jx_j\}.
\end{equation}
Let $\tau_n=\min\{j:N_j=0\}$.  Iterating
Proposition~\ref{prop:reciprocity} gives the exact identity
\begin{equation}\label{eq:exact-sum}
 S_n(\alpha)=\sum_{j=0}^{\tau_n-1}(-1)^j\Phi(x_j,u_j).
\end{equation}

Write
\[
 \alpha=[0;a_1,a_2,\ldots],\qquad x_j=T^j\alpha,
 \qquad a_{j+1}=a_1(x_j),
\]
and let $p_j/q_j$ be the convergents, normalized by
\[
 p_{-1}=1,\quad p_0=0,\qquad q_{-1}=0,\quad q_0=1.
\]
The standard identities
\begin{equation}\label{eq:beta}
 \beta_j:=x_0x_1\cdots x_j
 =(-1)^j(q_j\alpha-p_j)
 =\frac1{q_{j+1}+q_jx_{j+1}}
\end{equation}
will be used repeatedly; set $\beta_{-1}=1$.

Define the continuous heights and their errors by
\[
 C_j=n\beta_{j-1},\qquad E_j=C_j-N_j.
\]
Then $C_{j+1}=x_jC_j$, and \eqref{eq:descent} yields
\begin{equation}\label{eq:error-recurrence}
 E_{j+1}=x_jE_j+u_j,\qquad E_0=0.
\end{equation}
For every $r,\ell\ge0$,
\begin{equation}\label{eq:fibonacci-product}
 x_rx_{r+1}\cdots x_{r+\ell-1}\le F_{\ell+1}^{-1},
\end{equation}
where $F_1=F_2=1$.  Indeed, the reciprocal of the product is a
continuant with all digits at least one.  Iterating \eqref{eq:error-recurrence}
therefore gives
\begin{equation}\label{eq:E-bound}
 0\le E_j\le E_*:=\sum_{\ell\ge1}F_\ell^{-1}<\infty
\end{equation}
uniformly in $j,n,\alpha$.

Put
\begin{equation}\label{eq:theta}
 \theta_j=\{n\beta_j\}=\{(-1)^jnq_j\alpha\},
 \qquad \theta_{-1}=0.
\end{equation}
From $q_{j+1}=a_{j+1}q_j+q_{j-1}$ we obtain
\begin{equation}\label{eq:theta-cocycle}
 \theta_{j+1}=\{\theta_{j-1}-a_{j+1}\theta_j\}.
\end{equation}
Since $\{E_j\}=\theta_{j-1}$, write
\[
 E_j=d_j+\theta_{j-1},\qquad d_j=\lfloor E_j\rfloor.
\]
The equality $C_{j+1}=N_{j+1}+u_j+x_jE_j$ implies
\begin{equation}\label{eq:u-carry}
 u_j=\{\theta_j-x_j(d_j+\theta_{j-1})\}.
\end{equation}

Let
\[
 W(t)=\frac{\{t\}(1-\{t\})}{2},\qquad t\in\R/\Z.
\]
This function is $1/2$-Lipschitz and $0\le W\le1/8$.

\begin{proposition}[Principal term]\label{prop:principal}
There is an absolute constant $C_0$ such that, for every $j$,
\begin{equation}\label{eq:principal-split}
 \Phi(x_j,u_j)=a_{j+1}W(\theta_j)+B_j,
 \qquad |B_j|\le C_0.
\end{equation}
\end{proposition}

\begin{proof}
By \eqref{eq:u-carry} and \eqref{eq:E-bound},
\[
 \dist_{\T}(u_j,\theta_j)\le x_jE_*.
\]
Since $x_j^{-1}=a_{j+1}+x_{j+1}$,
\[
 \Phi(x_j,u_j)-a_{j+1}W(\theta_j)
 =a_{j+1}\{W(u_j)-W(\theta_j)\}
  +x_{j+1}W(u_j)-\frac{u_j}{2}.
\]
The first term has absolute value at most
$a_{j+1}x_jE_*/2\le E_*/2$, and the other two are bounded.
\end{proof}


\section{Gauss estimates and local oscillation}

The Gauss invariant probability measure is
\begin{equation}\label{eq:gauss-measure}
 d\nu(x)=h(x)\,dx,
 \qquad h(x)=\frac1{\log2\,(1+x)}.
\end{equation}
Its Lyapunov exponent is
\begin{equation}\label{eq:lambda}
 \lambda=\int_0^1-\log x\,d\nu(x)
 =\frac{\pi^2}{12\log2}.
\end{equation}

For a word $w=(a_1,\ldots,a_r)$ let $I_w$ be the corresponding
continued-fraction cylinder.  If $p_r/q_r$ is its last convergent, the inverse
branch is
\[
 h_w(y)=\frac{p_r+p_{r-1}y}{q_r+q_{r-1}y},
 \qquad
 |h_w'(y)|=(q_r+q_{r-1}y)^{-2},
\]
and
\begin{equation}\label{eq:cylinder-length}
 |I_w|=\frac1{q_r(q_r+q_{r-1})}.
\end{equation}
Conditionally on $I_w$, the density of $T^r\alpha=y$ is
\[
 \rho_w(y)=\frac{q_r(q_r+q_{r-1})}{(q_r+q_{r-1}y)^2}.
\]
Thus $\sup_w\|\rho_w\|_{BV}<\infty$.

We use the following standard consequences of the Lasota--Yorke inequality
for the Gauss transfer operator; see, for example, \cite{Baladi00}.  They are
recorded with the uniformity needed below.

\begin{lemma}[Gauss estimates]\label{lem:gauss-estimates}
There are $C>0$ and $0<\rho<1$ such that:
\begin{enumerate}
\item for $f\in BV(0,1)$,
\[
 \|\mathcal L^rf-h\!\int f\|_1\le C\rho^r\|f\|_{BV};
\]
\item conditionally on any past cylinder,
\[
 \Pp(a_{j+1}\ge A\mid a_1,\ldots,a_j)\le C/A;
\]
in particular, for distinct $j_1<\cdots<j_s$,
\begin{equation}\label{eq:multi-digit-tail}
 \Pp(a_{j_i+1}\ge A_i,\ 1\le i\le s)
 \le C^s\prod_{i=1}^s A_i^{-1};
\end{equation}
\item uniformly for $r\ge1$ and $v>0$,
\begin{equation}\label{eq:q-large-deviation}
 \Pp(|\log q_r-\lambda r|>v)
 \le C\exp\{-c\min(v^2/r,v)\}.
\end{equation}
\end{enumerate}
\end{lemma}

\begin{proof}
For completeness, the transfer operator is
\[
 \mathcal Lf(x)=\sum_{a\ge1}(a+x)^{-2}
 f\!\left((a+x)^{-1}\right).
\]
On two-step inverse branches the contraction is at most $1/2$; the sums of
the inverse Jacobians and of their variations converge.  The branchwise chain
rule gives
\[
 \Var(\mathcal L^2f)\le\rho_0\Var(f)+C\|f\|_1,
 \qquad \rho_0<1.
\]
Helly compactness and exactness yield (i).  The conditional densities
$\rho_w$ have uniformly bounded variation, and
$\{a_1\ge A\}=(0,1/A]$ up to endpoints; this proves (ii), and iteration proves
\eqref{eq:multi-digit-tail}.  Finally, the perturbed operators
\[
 \mathcal L_t f(x)=\sum_{a\ge1}(a+x)^{-2+t}
 f\!\left((a+x)^{-1}\right)
\]
form an analytic family on $BV$ for $|t|<1/2$.  Analytic perturbation of the
simple leading eigenvalue gives an eigenvalue $\kappa(t)$ and a uniformly
bounded spectral projector such that
\[
 \log\kappa(t)=\lambda t+O(t^2),
 \qquad
 \E\exp\!\left(t\sum_{i<r}-\log T^i\alpha\right)
 \le C\kappa(t)^r
\]
for $|t|$ sufficiently small; the same estimate holds for negative $t$.
For $0<v\le r$, Chernoff's inequality with $|t|\asymp v/r$ gives
$C\exp(-cv^2/r)$, while for $v>r$ a fixed sufficiently small $|t|$ gives
$C\exp(-cv)$.  Thus the asserted Bernstein bound holds for
$\sum_{i<r}-\log T^i\alpha$.
Since
\[
 \left|\log q_r-\sum_{i<r}-\log T^i\alpha\right|\le\log2,
\]
we obtain (iii).
\end{proof}

\begin{lemma}[Conditional multi-block mixing]\label{lem:conditional-mixing}
Let $I_w$ be a cylinder of depth $d$, and let
$d+M\le t_1<\cdots<t_s$ with $t_{i+1}-t_i\ge M$.  If
$g_1,\ldots,g_s\in BV(0,1)$, then
\begin{equation}\label{eq:conditional-mixing}
 \left|
 \E\left(\prod_{i=1}^s g_i(T^{t_i}\alpha)\,\middle|\,I_w\right)
 -\prod_{i=1}^s\int g_i\,d\nu
 \right|
 \le C_s\rho^M\prod_{i=1}^s\|g_i\|_{BV}.
\end{equation}
The constants are uniform in the cylinder $I_w$.
\end{lemma}

\begin{proof}
Conditioned on $I_w$, the density of $T^d\alpha$ is $\rho_w$, whose $BV$
norm is uniformly bounded by \eqref{eq:cylinder-length}.  Apply
Lemma~\ref{lem:gauss-estimates}(i) across the first gap, multiply by the next
$BV$ observable, and iterate.  The product inequality
$\|fg\|_{BV}\le C\|f\|_{BV}\|g\|_{BV}$ and the fixed number $s$ of factors
give \eqref{eq:conditional-mixing}.
\end{proof}

Fix henceforth
\begin{equation}\label{eq:scales}
 L=\log n,\qquad H=L^{3/4},\qquad
 m_n=\left\lfloor L/\lambda\right\rfloor,
\end{equation}
and define the bulk set
\begin{equation}\label{eq:bulk}
 \mathcal J_n=\{j\in\Z:200H\le j\le m_n-200H\}.
\end{equation}
By \eqref{eq:q-large-deviation}, for every fixed $\delta>0$ and
$t\le2m_n$,
\begin{equation}\label{eq:local-q-good}
 \Pp(q_t\notin[\e^{\lambda t-\delta H},
                    \e^{\lambda t+\delta H}])
 \le C_\delta \e^{-c_\delta L^{1/2}}.
\end{equation}
We always use \eqref{eq:local-q-good} as a union of complete depth-$t$
cylinders.  We never insert a global good-event indicator inside an
oscillatory integral.

The first ingredient controls local cancellation of the two continuants in a
torus character.

\begin{lemma}[Shrinking anti-concentration]\label{lem:anti}
For every fixed $(r,s)\in\Z^2\setminus\{(0,0)\}$, every $j\ge1$, and
$0<\eta<1/2$,
\begin{equation}\label{eq:anti}
 \Pp(|sq_j-rq_{j-1}|<\eta q_j)
 \le C_{r,s}\eta+C_{r,s}\e^{-cj}.
\end{equation}
\end{lemma}

\begin{proof}
Put $Y_j=q_{j-1}/q_j=[0;a_j,\ldots,a_1]$ and
$Y_{j,R}=[0;a_j,\ldots,a_{j-R+1}]$.  The contraction of inverse continued
fractions gives
\[
 |Y_j-Y_{j,R}|\le C\varphi^{-2R},
 \qquad \varphi=(1+\sqrt5)/2.
\]
The Gauss natural extension has invariant density
\[
 \frac{dx\,dy}{\log2\,(1+xy)^2};
\]
its backward marginal is $dy/(\log2(1+y))$.  By
Lemma~\ref{lem:gauss-estimates}(i), the law of the last $R$ digits under
Lebesgue measure differs in total variation from its stationary law by
$O(\rho^{j-R})$.  Hence, for every interval $I\subset(0,1)$,
\[
 \Pp(Y_j\in I)\le C(|I|+\varphi^{-2R})+C\rho^{j-R}.
\]
Taking $R=\lfloor j/2\rfloor$ gives $C|I|+C\e^{-cj}$.  If $r\ne0$, the event
in \eqref{eq:anti} is $|s-rY_j|<\eta$ and is contained in an interval of
length $2\eta/|r|$.  If $r=0$, it is empty.
\end{proof}

The second ingredient is a complete-cylinder oscillatory estimate.  Its
formulation makes the measurability needed for integration by parts explicit.

\begin{lemma}[Descendant-cylinder estimate]\label{lem:descendant}
Let $\{I_w\}$ be cylinders of a fixed depth $d$.  On $I_w$ let
$Q_w\in\Z\setminus\{0\}$ be fixed.  Let $k>d$ and let $A$ be constant on every
depth-$k$ cylinder, with $|A|\le1$.  Suppose that, below $I_w$, the support of
$A$ is a union of depth-$k$ descendants satisfying $q_k\le R_w$, where
\[
 R_w^2\le\varepsilon n|Q_w|.
\]
Then
\begin{equation}\label{eq:descendant}
 \left|\sum_w\int_{I_w}
       \e^{2\pi i nQ_w\alpha}A(\alpha)\,d\alpha\right|
 \le C\varepsilon.
\end{equation}
Conversely, if a depth-$k$ cylinder $J$ satisfies
$q_k^2\ge\varepsilon^{-1}n|Q|$, then
\begin{equation}\label{eq:phase-freeze}
 \sup_{\alpha,\alpha'\in J}|nQ(\alpha-\alpha')|
 \le C\varepsilon.
\end{equation}
\end{lemma}

\begin{proof}
Let $q(w)$ be the last denominator of $w$.  If $(Q,Q')$ is the denominator
pair of a continuation, the final denominator is
$q_k=q(w)Q+q'(w)Q'\ge q(w)Q$.  The number of descendants with
$q_k\le R_w$ is therefore $O(R_w^2/q(w)^2)$.  On each complete descendant
cylinder $J$,
\[
 \left|\int_J\e^{2\pi i nQ_w\alpha}\,d\alpha\right|
 \le(\pi n|Q_w|)^{-1}.
\]
Thus the contribution below $w$ is at most $C\varepsilon/q(w)^2$.  By
\eqref{eq:cylinder-length}, $q(w)^{-2}\le2|I_w|$, and summation proves
\eqref{eq:descendant}.  The cylinder length bound $|J|\le q_k^{-2}$ gives
\eqref{eq:phase-freeze}.
\end{proof}


\section{Marked factorization and resonances}

For $D>0$ let $\mathcal P_D(L)$ consist of functions
\begin{equation}\label{eq:one-time-poly}
 F(a,\theta)=\sum_{|v|\le L^D}c(a,v)\e^{2\pi iv\theta},
 \qquad c(a,v)=0\quad(a>L^D),
\end{equation}
such that $\sum_{a,v}|c(a,v)|\le L^D$.  For times
$j_1<\cdots<j_r$ in $\mathcal J_n$, call the tuple good if
\begin{equation}\label{eq:good-tuple-separation}
 j_{\ell+1}-j_\ell\ge200H
\end{equation}
and, for every $i<\ell$,
\begin{equation}\label{eq:good-tuple-resonance}
 \left|j_\ell-\frac{m_n+j_i}{2}\right|\ge200H.
\end{equation}
The number of non-good $r$-tuples is $O_r(L^{r-1}H)$.

\begin{proposition}[Marked factorization]\label{prop:marked-factorization}
Fix $r,D,A$.  If $(j_1,\ldots,j_r)$ is good and
$F_1,\ldots,F_r\in\mathcal P_D(L)$, then
\begin{align}\label{eq:marked-factorization}
 &\int_0^1\prod_{\ell=1}^r
 F_\ell(a_{j_\ell+1},\theta_{j_\ell})\,d\alpha\\
 &\qquad=
 \prod_{\ell=1}^r\int_0^1\int_0^1
 F_\ell(a_1(x),\theta)\,d\theta\,d\nu(x)
 +O_{r,D,A}(L^{-A}).\notag
\end{align}
\end{proposition}

\begin{proof}
Expand in Fourier modes $v_1,\ldots,v_r$.  If every $v_\ell=0$, the
integrand is a product of digit functions.  Repeated use of
Lemma~\ref{lem:conditional-mixing} and
\eqref{eq:good-tuple-separation} factors its mean with error
$L^{O_{r,D}(1)}\rho^{200H}=O_A(L^{-A})$.

Suppose now that $v_s\ne0$ and $v_\ell=0$ for $\ell>s$.  By
\eqref{eq:theta}, the phase is
\[
 \e^{2\pi i nQ\alpha},
 \qquad
 Q=\sum_{\ell=1}^{s}v_\ell(-1)^{j_\ell}q_{j_\ell}.
\]
The continuants grow at least at the Fibonacci rate.  The separation
condition and $|v_\ell|\le L^D$ therefore give, deterministically,
\begin{equation}\label{eq:Q-dominance}
 \tfrac12q_{j_s}\le|Q|\le2L^Dq_{j_s}
\end{equation}
for all sufficiently large $n$.

The nominal resonance time is
\[
 R_s=\frac{m_n+j_s}{2}.
\]
Put $k_\pm=\lfloor R_s\pm40H\rfloor$.  We now make only local complete-cylinder
cuts.  At depth $j_s$ retain the cylinder union on which
\[
 \e^{\lambda j_s-\delta H}\le q_{j_s}
 \le \e^{\lambda j_s+\delta H}.
\]
At depths $k_-$ and $k_+$ retain, respectively,
\[
 q_{k_-}\le \e^{\lambda k_-+\delta H},\qquad
 q_{k_+}\ge \e^{\lambda k_+-\delta H},
\]
where $\delta>0$ is sufficiently small.  By
\eqref{eq:local-q-good}, the total discarded mass is
$O(\e^{-cL^{1/2}})$.  On every retained prefix and descendant,
\begin{equation}\label{eq:pre-post-resonance}
 q_{k_-}^2\le \e^{-cH}n|Q|,
 \qquad
 q_{k_+}^2\ge \e^{cH}n|Q|.
\end{equation}

The factors after $s$ are zero torus modes and hence digit functions.  By
\eqref{eq:good-tuple-resonance}, each lies either at least $100H$ before
$k_-$ or at least $100H$ after $k_+$.  Partition first into complete
depth-$k_+$ cylinders.  On each retained cylinder the phase is constant up to
$O(\e^{-cH})$ by \eqref{eq:phase-freeze}.  The conditional density of the
future has uniformly bounded variation, so
Lemma~\ref{lem:conditional-mixing}
replaces all post-resonance digit factors by their stationary means, with
total error $L^{O(1)}(\e^{-cH}+\rho^{cH})$ after summation by cylinder mass.

After that replacement, extend the stationary-mean replacement to the
discarded depth-$k_+$ cylinders and restore them.  Their total mass is
$O(\e^{-cL^{1/2}})$, so this costs at most
$L^{O_{r,D}(1)}\e^{-cL^{1/2}}$.  The depth-$k_+$ cutoff indicator has now
disappeared, and the remaining amplitude is measurable on complete
depth-$k_-$ descendants.  On each depth-$(j_s+1)$ prefix the integer $Q$ is
fixed.  The first inequality in \eqref{eq:pre-post-resonance} and
Lemma~\ref{lem:descendant} give a contribution $O(\e^{-cH})$.  Restoring the
discarded cylinder unions, and summing all polynomially many Fourier terms,
gives the uniform bound
\begin{equation}\label{eq:delta-n}
 \delta_n:=L^{O_{r,D}(1)}
 \bigl(\e^{-cL^{1/2}}+\e^{-cH}+\rho^{cH}\bigr)=O_A(L^{-A}).
\end{equation}
Thus every nonzero Fourier term vanishes to the required order, while the
zero terms give the product on the right of
\eqref{eq:marked-factorization}.
\end{proof}

We also need fixed-radius blocks.  The recurrence
\eqref{eq:theta-cocycle} implies by induction that every torus character in
a radius-$R$ window can be reduced, on its digit cylinder, to a character of
the two central coordinates.  More precisely, for fixed integers
$c_{-R},\ldots,c_R$ there are integers $A_w,B_w$, depending only on the
radius-$R$ digit word $w$, such that
\begin{equation}\label{eq:block-character-reduction}
 \sum_{|t|\le R}c_t\theta_{j+t}
 \equiv A_w\theta_j+B_w\theta_{j-1}\pmod1.
\end{equation}
This follows by repeatedly substituting
$\theta_{i+1}=\theta_{i-1}-a_{i+1}\theta_i$.  On each fixed finite
collection of radius-$R$ cylinder words, $A_w,B_w$ range over a finite set.
In view of
\eqref{eq:theta}, the resulting phase is
$n(-1)^j(A_wq_j-B_wq_{j-1})\alpha$.

Let $R,K$ be fixed and call a cylinder--torus monomial a function of the form
\begin{equation}\label{eq:block-monomial}
 U_j(\alpha)=u(a_{j-R+1},\ldots,a_{j+R})
 \e^{2\pi i(r\theta_{j-1}+s\theta_j)},
\end{equation}
where $u$ is supported on finitely many words and $|r|+|s|\le K$.  Its
frequency is
\begin{equation}\label{eq:block-frequency}
 (-1)^jQ_j(r,s),\qquad Q_j(r,s)=sq_j-rq_{j-1}.
\end{equation}

\begin{proposition}[Two-block factorization]\label{prop:two-block}
For two fixed finite linear combinations $U,V$ of monomials
\eqref{eq:block-monomial}, all but $O_{U,V}(LH)$ pairs
$j<k$ in $\mathcal J_n$ satisfy
\begin{equation}\label{eq:two-block}
 \left|\int_0^1U_jV_k\,d\alpha
 -\int U\,d\widehat\mu_0\int V\,d\widehat\mu_0\right|
 \le C_{U,V}\bigl(\e^{-cL^{1/2}}+\e^{-cH}+\rho^{cH}\bigr),
\end{equation}
where $\widehat\mu_0$ is the stationary Gauss natural-extension measure
times Haar measure on $\T^2$.
\end{proposition}

\begin{proof}
It suffices to treat monomials.  If both torus modes vanish, ordinary
two-sided Gauss mixing proves \eqref{eq:two-block} unless the blocks overlap;
there are $O(LH)$ overlapping or $H$-near pairs.

Assume the later mode $(r_2,s_2)$ is nonzero.  By
Lemma~\ref{lem:anti}, outside a complete depth-$k$ cylinder union of mass
$O(\e^{-cH})$,
\[
 |Q_k(r_2,s_2)|\ge \e^{-cH}q_k.
\]
If $k-j\ge CH$, Fibonacci growth makes the later frequency dominate the
earlier one; after increasing $C=C(U,V)$,
\[
 |Q|:=|(-1)^jQ_j(r_1,s_1)+(-1)^kQ_k(r_2,s_2)|
 \ge\tfrac12\e^{-cH}q_k.
\]
On each depth-$(k+R)$ cylinder, the integer $Q$ and the complete amplitude
are fixed.  Put
\[
 t_-:=\left\lfloor\frac{m_n+k}{2}-40H\right\rfloor.
\]
Retain the complete depth-$k$ cylinders on which
$q_k\ge \e^{\lambda k-\delta H}$ and, below each such cylinder, the complete
depth-$t_-$ descendants on which
$q_{t_-}\le \e^{\lambda t_-+\delta H}$.  By
\eqref{eq:local-q-good}, the discarded union has mass
$O(\e^{-cL^{1/2}})$.  Since
$2\lambda t_-=L+\lambda k-80\lambda H+O(1)$, every retained descendant
satisfies
\[
 q_{t_-}^2\le \e^{-cH}n|Q|.
\]
Lemma~\ref{lem:descendant}, with prefix depth $k+R$ and descendant depth
$t_-$, bounds the retained integral by $O(\e^{-cH})$.  The summation is by
complete prefix-cylinder mass, so no cylinder-count factor is lost.  Adding
the anti-concentration and denominator-deviation unions gives the right side
of \eqref{eq:two-block}.  The stationary product is zero because the later
Haar character is nontrivial.

It remains to suppose that the earlier mode is nonzero and the later mode is
zero.  Write $Q_j=Q_j(r_1,s_1)$ and
\[
 t_0:=\frac{m_n+j}{2},\qquad
 t_-:=\lfloor t_0-40H\rfloor,
 \qquad t_+:=\lfloor t_0+40H\rfloor.
\]
Lemma~\ref{lem:anti} removes a complete depth-$j$ cylinder union of mass
$O(\e^{-cH})$ and leaves $|Q_j|\ge\e^{-cH}q_j$.  Declare
$|k-t_0|\le100H$ exceptional.  This gives $O(H)$ values of $k$ per $j$,
hence $O(LH)$ pairs in total.

We also retain throughout the complete depth-$j$ cylinder union on which
\[
 \e^{\lambda j-\delta H}\le q_j\le\e^{\lambda j+\delta H}.
\]
By \eqref{eq:local-q-good}, its complement has mass
$O(\e^{-cL^{1/2}})$.  On the retained cylinders,
\[
 \e^{-cH}q_j\le |Q_j|\le C_{r_1,s_1}q_j.
\]

If $k<t_0-100H$, the whole two-block amplitude is measurable before depth
$t_-$.  Under the retained depth-$j$ prefixes, keep the complete
depth-$t_-$ descendants with $q_{t_-}\le\e^{\lambda t_-+\delta H}$.
The complement has mass
$O(\e^{-cL^{1/2}})$, while on the retained descendants
\[
 q_{t_-}^2\le\e^{-cH}n|Q_j|.
\]
Lemma~\ref{lem:descendant} removes the phase with error $O(\e^{-cH})$.

If $k>t_0+100H$, also retain the complete depth-$t_+$ cylinders satisfying
$q_{t_+}\ge\e^{\lambda t_+-\delta H}$.  Their complement again has mass
$O(\e^{-cL^{1/2}})$ and, on every retained cylinder,
\[
 q_{t_+}^2\ge\e^{cH}n|Q_j|.
\]
Thus \eqref{eq:phase-freeze} makes the earlier phase constant up to
$O(\e^{-cH})$ on each retained depth-$t_+$ cylinder.  The later zero-mode
block starts at least $cH$ after $t_+$, so conditional Gauss mixing, applied
on each cylinder and summed by cylinder mass, replaces that block by its
stationary mean with error $O(\rho^{cH})$.  Extend this stationary-mean
replacement to the discarded depth-$t_+$ cylinders and restore them, at cost
$O(\e^{-cL^{1/2}})$; no future cutoff indicator then remains in the earlier
integral.  Finally retain the depth-$t_-$ descendants
as above and apply Lemma~\ref{lem:descendant} to the remaining earlier phase.

The stationary product is zero because the earlier Haar character is
nontrivial.  Combining the complete-cylinder errors proves
\eqref{eq:two-block}; in particular, the denominator-deviation contribution
is $O(\e^{-cL^{1/2}})$ rather than being hidden in the anti-concentration
error.
\end{proof}



\section{The principal Poisson and stable limits}

Let
\[
 Z_{n,j}=a_{j+1}W(\theta_j).
\]
If $A$ has the stationary Gauss digit law and $\Theta$ is independent and
uniform on $\T$, then
\begin{equation}\label{eq:stationary-tail}
 \Pp(AW(\Theta)>t)\sim\frac{c_0}{t},
 \qquad c_0=\frac1{12\log2}.
\end{equation}
Indeed,
\[
 \nu(A\ge m)=\frac1{\log2}\log\left(1+\frac1m\right)
 \sim\frac1{m\log2},
\]
and dominated convergence gives
\[
 \lim_{t\to\infty}t\Pp(AW(\Theta)>t)
 =\frac1{\log2}\int_0^1W(\theta)\,d\theta
 =\frac1{12\log2}.
\]

\begin{proposition}[Poisson point process]\label{prop:PPP}
In the space of Radon point measures on $\R\setminus\{0\}$ equipped with
the vague topology,
\begin{equation}\label{eq:PPP}
 \sum_{j\in\mathcal J_n}\delta_{(-1)^jZ_{n,j}/L}
 \ \Longrightarrow\
 \operatorname{PRM}(\Lambda),
 \qquad
 d\Lambda(x)=\frac1{2\pi^2}|x|^{-2}\,dx.
\end{equation}
\end{proposition}

\begin{proof}
We verify factorial moment measures against
$\varphi_1,\ldots,\varphi_r\in C_c^1(\R\setminus\{0\})$.
Writing $X_{n,j}=(-1)^jZ_{n,j}/L$, the relevant ordered factorial sum is
\begin{equation}\label{eq:ordered-factorial-sum}
 \sum_{\substack{(j_1,\ldots,j_r)\in\mathcal J_n^r\\
                  j_1,\ldots,j_r\ \text{ pairwise distinct}}}
 \E\prod_{\ell=1}^r\varphi_\ell(X_{n,j_\ell}).
\end{equation}
Let $c>0$ be smaller than the distance of all supports from zero.  If
$\varphi_\ell((-1)^jaW(\theta)/L)\ne0$, then $a\ge8cL$.

First cut at $a\le L^D$, where $D>2$.  By
\eqref{eq:multi-digit-tail}, the total contribution of $r$-tuples for which
one digit exceeds $L^D$ and all the others exceed $8cL$ is
\[
 O\left(L^r\frac1{L^D L^{r-1}}\right)=O(L^{1-D}).
\]
For $a\le L^D$, the $C^1$ function
$f_{a,j}(\theta)=\varphi_\ell((-1)^jaW(\theta)/L)$ satisfies
$\|f_{a,j}'\|_\infty\le C_\varphi a/L\le C_\varphi L^{D-1}$.
Jackson's theorem \cite[Chap.~X]{Zygmund59}, with degree
$M=L^{A+r+D+4}$, gives a trigonometric
polynomial whose uniform error is $O(L^{-A-r-4})$.  Its coefficient
$\ell^1$ norm is at most
$(2M+1)(\|\varphi_\ell\|_\infty+1)$ and hence is polynomial in $L$.
Increasing the complexity exponent in
Proposition~\ref{prop:marked-factorization}, we may apply that proposition;
after summation over $O(L^r)$ tuples, the approximation error is
$O(L^{-A-4})$.

Non-good time tuples are $O_r(L^{r-1}H)$ in number.  On each of them the
probability that all factors are nonzero is $O_r(L^{-r})$ by
\eqref{eq:multi-digit-tail}; their total contribution is $O(H/L)=o(1)$.
On good tuples, Proposition~\ref{prop:marked-factorization} factors the
expectation, and its error remains $o(1)$ after summing $O(L^r)$ tuples.
Equation \eqref{eq:stationary-tail} gives the vague convergence
\[
 L\Pp(AW(\Theta)/L\in dx)\Longrightarrow c_0x^{-2}\,dx
 \quad(x>0).
\]
To evaluate \eqref{eq:ordered-factorial-sum}, partition the ordered distinct
tuples according to the unique permutation $\sigma\in S_r$ for which
$j_{\sigma(1)}<\cdots<j_{\sigma(r)}$, and then according to the parity vector
$\varepsilon_\ell=(-1)^{j_\ell}\in\{-1,1\}$.  For every fixed
$\sigma$ and $\varepsilon=(\varepsilon_1,\ldots,\varepsilon_r)$,
\begin{equation}\label{eq:parity-tuple-count}
 \#\left\{(j_1,\ldots,j_r):
 j_{\sigma(1)}<\cdots<j_{\sigma(r)},\
 (-1)^{j_\ell}=\varepsilon_\ell\right\}
 =\frac{1+o(1)}{r!}\left(\frac{L}{2\lambda}\right)^r.
\end{equation}
Indeed, the parity-restricted lattice discrepancy in the actual bulk interval
is $O(L^{r-1})$, while replacing its length by $L/\lambda$ contributes
$O(HL^{r-1})=o(L^r)$.  Put
\[
 I_\ell^{\varepsilon}
 =\int_0^\infty\varphi_\ell(\varepsilon x)x^{-2}\,dx.
\]
For a fixed ordering and parity vector, factorization and
\eqref{eq:stationary-tail} show that the contribution tends to
\[
 \frac1{r!}\left(\frac{c_0}{2\lambda}\right)^r
 \prod_{\ell=1}^r I_\ell^{\varepsilon_\ell}.
\]
Summing first over the $r!$ temporal orderings and then over all parity
vectors yields
\[
 \prod_{\ell=1}^r
 \left\{\frac{c_0}{2\lambda}
 \int_{\R}\varphi_\ell(x)|x|^{-2}\,dx\right\}
 =\prod_{\ell=1}^r\int\varphi_\ell\,d\Lambda,
 \qquad \frac{c_0}{2\lambda}=\frac1{2\pi^2}.
\]
Thus the ordered factorial moment measures converge to those of the Poisson
random measure in \eqref{eq:PPP}.  If $K\Subset\R\setminus\{0\}$ and
$N_n(K)$ denotes the number of displayed points in $K$, the same iterated
conditional tail bound gives, for every $r\ge1$,
\[
 \sup_n\E (N_n(K))_r\le C_K^r,
\]
where $(x)_r=x(x-1)\cdots(x-r+1)$.  In particular all local factorial
moments are uniformly bounded, while the first moment supplies local
tightness.  The
factorial-moment criterion for point processes
\cite[Chap.~4]{Kallenberg17}, applied on an exhausting sequence of compact
continuity sets, identifies every subsequential limit as the Poisson random
measure with intensity $\Lambda$.
\end{proof}

We next control the compensated small jumps.  Choose
$\chi\in C^1([0,\infty))$ with $0\le\chi\le1$, $\chi=1$ on $[0,1]$, and
$\chi=0$ on $[2,\infty)$.  For $0<\varepsilon<1$, set
\begin{equation}\label{eq:small-jump-U}
 U_{n,j}^{(\varepsilon)}=(-1)^jZ_{n,j}
 \chi\!\left(\frac{Z_{n,j}}{\varepsilon L}\right).
\end{equation}

\begin{lemma}[Small jumps]\label{lem:small-jumps}
For a constant $C$ independent of $\varepsilon$,
\begin{equation}\label{eq:small-jump-var}
 \limsup_{n\to\infty}\frac1{L^2}
 \Var\left(\sum_{j\in\mathcal J_n}U_{n,j}^{(\varepsilon)}\right)
 \le C\varepsilon.
\end{equation}
\end{lemma}

\begin{proof}
Since $W\le1/8$, Lemma~\ref{lem:gauss-estimates}(ii) gives, uniformly in
$j,n$,
\[
 \Pp(Z_{n,j}>t)\le\frac{C}{1+t}.
\]
For $j\ne k$, \eqref{eq:multi-digit-tail} similarly gives
\[
 \Pp(Z_{n,j}>s,Z_{n,k}>t)
 \le\frac{C}{(1+s)(1+t)}.
\]
Layer-cake integration, with $M=2\varepsilon L$, yields
\begin{equation}\label{eq:small-moments}
 \E|U_{n,j}^{(\varepsilon)}|^2\le C\varepsilon L,
 \qquad
 \E|U_{n,j}^{(\varepsilon)}U_{n,k}^{(\varepsilon)}|
 \le C\log^2(2+L).
\end{equation}
The diagonal contribution is $O(\varepsilon L^2)$.  Overlapping,
$H$-near, or resonance-near pairs number $O(LH)$, and their total
contribution is $O(LH\log^2L)=o(L^2)$.

For the remaining pairs, cut each digit at $a\le A_L=L^D$.  Since
$|U_{n,j}^{(\varepsilon)}|\le2\varepsilon L$,
\[
 \|U_{n,j}^{(\varepsilon)}\ind_{\{a_{j+1}>A_L\}}\|_2
 \le C\varepsilon L A_L^{-1/2}.
\]
After centering, $\|X-\E X\|_2\le2\|X\|_2$; hence, after division by $L$,
Minkowski's inequality shows that the sum of these tails is
$O(\varepsilon L^{1-D/2})=o(1)$ if $D>2$.  On $a\le A_L$, the
derivative of the truncated summand is polynomial in $L$; Jackson
approximation with uniform error $L^{-A-4}$ and polynomial degree and
coefficient norm applies as in the proof of Proposition~\ref{prop:PPP}.
Centering doubles the approximation error but does not change its order.
Proposition~\ref{prop:marked-factorization},
with $A$ chosen so large that $L^2O_A(L^{-A})=o(1)$, makes the sum of all
remaining covariances $o(L^2)$.  This proves
\eqref{eq:small-jump-var}.
\end{proof}

\begin{corollary}[Principal Cauchy law]\label{cor:principal-cauchy}
There are deterministic numbers $b_n^{(0)}$ such that
\begin{equation}\label{eq:principal-cauchy}
 \frac1L\sum_{j\in\mathcal J_n}(-1)^jZ_{n,j}-b_n^{(0)}
 \Longrightarrow \operatorname{Cauchy}\!\left(0,\frac1{2\pi}\right).
\end{equation}
\end{corollary}

\begin{proof}
For fixed $\varepsilon$, split $x=x\chi(|x|/\varepsilon)
+x(1-\chi(|x|/\varepsilon))$.  Proposition~\ref{prop:PPP} and the
continuous mapping theorem give convergence of the second, large-jump part;
the digit-tail bound gives the uniform finite-$n$ estimate
\[
 \sup_n\E\Xi_n\{x:|x|>R\}\le C/R,
\]
so the tails $|x|>R$ may be removed before applying continuous mapping.
Center the small part by
\[
 b_{n,\varepsilon}=\frac1L\sum_{j\in\mathcal J_n}
 \E U_{n,j}^{(\varepsilon)}.
\]
Lemma~\ref{lem:small-jumps} shows that its centered contribution becomes
negligible as $\varepsilon\downarrow0$.  The transition shell
$\varepsilon<|x|<2\varepsilon$ is included in the same continuous
truncation; its limiting compensated variance is $O(\varepsilon)$.

For the smooth cutoff used here, the limiting large-jump characteristic
exponent is, first on $|x|\le R$,
\[
 \int_{|x|\le R}
 \left(\exp\{itx(1-\chi(|x|/\varepsilon))\}-1\right)d\Lambda(x).
\]
Letting $R\to\infty$ is legitimate because the expected number of points
outside $[-R,R]$ is $O(R^{-1})$.  Symmetry cancels the imaginary part, and
then dominated convergence as $\varepsilon\downarrow0$ gives
\[
 \int_{\R}(\cos(tx)-1)\frac{dx}{2\pi^2x^2}
 =-\frac{|t|}{2\pi}.
\]
A diagonal choice $\varepsilon_n\downarrow0$ gives a deterministic sequence
$b_n^{(0)}:=b_{n,\varepsilon_n}$ for which
\eqref{eq:principal-cauchy} holds.
\end{proof}





\section{The carry graph and the bounded remainder}

Let $(\widehat\Omega,\widehat\nu,\sigma)$ be the two-sided natural extension
of the Gauss system and write $x(\omega)=[0;a_1(\omega),a_2(\omega),\ldots]$.
Define the torus cocycle
\begin{equation}\label{eq:S-cocycle}
 \mathcal S(\omega,r,s)=
 (\sigma\omega,s,r-a_1(\omega)s\pmod1).
\end{equation}
The fibre matrix
\[
 A_a=\begin{pmatrix}0&1\\1&-a\end{pmatrix}
\]
lies in $\operatorname{GL}_2(\Z)$; hence
$\widehat\mu_0=\widehat\nu\otimes m_{\T^2}$ is invariant.

\begin{lemma}[Endpoint-controlled continuant recurrence]
\label{lem:endpoint-recurrence}
Let $a_i\ge1$ and let a nonzero integer sequence satisfy
\[
 z_{i+2}=a_{i+1}z_{i+1}+z_i\qquad(0\le i<m).
\]
If
\[
 \max(|z_0|,|z_1|,|z_m|,|z_{m+1}|)\le K,
\]
then $m\le C\log(2K)$ for an absolute constant $C$.
\end{lemma}

\begin{proof}
Consider a maximal block on which consecutive nonzero terms have alternating
signs.  Write an interior pair as $(z_i,z_{i+1})=(u,-v)$ after changing the
overall sign, where $u,v>0$.  If alternation persists for two further terms,
then
\[
 z_{i+2}=u-a_{i+1}v>0,
 \qquad
 z_{i+3}=a_{i+2}z_{i+2}-v<0.
\]
Consequently $0<z_{i+2}<v$, and the absolute-value pair is replaced by the
ordinary Euclidean remainder pair $(v,u-a_{i+1}v)$.  Thus every interior
transition of the maximal alternating block is a Euclidean remainder step;
only its terminal transition can fail the remainder inequality, and that
costs one additional step.  Lam\'e's bound for the Euclidean
algorithm (whose extremal quotients are all $1$) bounds its length by
$C\log(2K)$ when both endpoint pairs are bounded by $K$.  A zero term can
occur only at the junction of two such intervals and causes no additional
run.

Once two adjacent nonzero terms have the same sign, every later adjacent pair
has that sign and
\[
 |z_{i+2}|\ge |z_{i+1}|+|z_i|;
\]
hence the forward absolute values grow at least as Fibonacci numbers.  The
same argument applied to the recurrence solved backwards controls the final
same-sign run.  The sequence is therefore the concatenation of at most two
Euclidean runs and one Fibonacci-growth run, each of length
$O(\log(2K))$.
\end{proof}

\begin{lemma}[Gauss--torus mixing]\label{lem:torus-mixing}
The system $(\widehat\Omega\times\T^2,\widehat\mu_0,\mathcal S)$ is mixing.
\end{lemma}

\begin{proof}
It is enough to use finite digit-cylinder functions times torus characters.
For characters $k,\ell\in\Z^2$, the fibre integral in an $m$-step
correlation vanishes unless
\begin{equation}\label{eq:character-resonance}
 k+(A_{a_m}\cdots A_{a_1})^{\mathrm T}\ell=0.
\end{equation}
After alternating signs are absorbed, either coordinate of the successive
integer vectors obeys the recurrence in
Lemma~\ref{lem:endpoint-recurrence}.  Thus, for fixed nonzero $k,\ell$, the resonance relation is
impossible for all sufficiently large $m$.  The zero modes reduce to mixing
of the Gauss natural extension.
Density of cylinder functions times characters in $L^2$ completes the proof.
\end{proof}

We need to transfer fixed windows of the actual curve process to this
stationary system.  The complete quotient is an infinite-future coordinate,
so we record the required full-state form rather than silently treating it as
a finite digit function.

For a fixed $R\ge0$, give
\begin{equation}\label{eq:window-space}
 \mathcal X_R
 =\N^{2R+1}\times[0,1]^{2R+1}\times\T^{2R+2}
\end{equation}
the product topology, with $\N$ discrete, and define the actual window
\begin{equation}\label{eq:actual-window}
 \mathcal W_{n,j}^{(R)}=
 \left((a_{j+t+1})_{t=-R}^{R},
       (x_{j+t})_{t=-R}^{R},
       (\theta_{j+t})_{t=-R-1}^{R}\right)\in\mathcal X_R.
\end{equation}
The stationary window is defined from the two-sided Gauss--torus system in
the same way.  Both laws are supported on the corresponding orbit-consistent
Borel subset of $\mathcal X_R$.

\begin{lemma}[Uniform full-state transfer]\label{lem:full-state}
Let $\eta_{n,j}^{(R)}$ be the law of \eqref{eq:actual-window} for Lebesgue
$\alpha$, and let $\mu_R$ be the stationary window law.  If
$G:\mathcal X_R\to\mathbb C$ is bounded and its discontinuity set has
$\mu_R$-measure zero, then
\begin{equation}\label{eq:full-state}
 \sup_{j\in\mathcal J_n}
 \left|\int G\,d\eta_{n,j}^{(R)}-\int G\,d\mu_R\right|\longrightarrow0.
\end{equation}
\end{lemma}

\begin{proof}
First take a finite digit-cylinder function times a torus character.  Partition
it into its finitely many complete depth-$(j+R)$ block cylinders.  By
\eqref{eq:block-character-reduction}, on each such cylinder the coefficients
$A_w,B_w$, the integer $Q_j=A_wq_j-B_wq_{j-1}$, and the whole block amplitude
are fixed.  If $A_w=B_w=0$, uniform Gauss mixing gives the stationary mean.
Otherwise Lemma~\ref{lem:anti}, with $\eta=\e^{-cH}$, removes, after summing
over the finite block collection, a complete cylinder union of mass
$O(\e^{-cH})$ on which local cancellation may occur.  On the complement,
$|Q_j|\ge\e^{-cH}q_j$.

Put $t=\lfloor(m_n+j)/2-40H\rfloor$.  Because
$j\in\mathcal J_n$, one has $t>j+R$.  Retain the coarser complete depth-$j$
cylinder union on which $q_j\ge\e^{\lambda j-\delta H}$ and, below every
retained depth-$(j+R)$ block prefix, the complete depth-$t$ descendants on
which $q_t\le\e^{\lambda t+\delta H}$.  The
discarded mass is $O(\e^{-cL^{1/2}})$ by
\eqref{eq:local-q-good}; on every retained descendant,
\[
 q_t^2\le\e^{-cH}n|Q_j|.
\]
Lemma~\ref{lem:descendant}, with prefix depth $j+R$, therefore makes every
nonzero character
$O(\e^{-cH}+\e^{-cL^{1/2}})$.  This proves the required one-block convergence
uniformly in $j$.  Finite sums give convergence on the cylinder--character
algebra.

Continuous functions of the complete quotients are included by replacing,
simultaneously for every $|t|\le R$, the coordinate
$x_{j+t}=[0;a_{j+t+1},a_{j+t+2},\ldots]$ by its first $M$ future digits.
The maximum coordinate error is $O_R(F_M^{-2})$; hence, for continuous $G$,
the test-function error is at most $\omega_G(C_RF_M^{-2})$.  Before invoking
the finite cylinder--character algebra, restrict all digits used in these
truncations to $a\le K$.  Lemma~\ref{lem:gauss-estimates}(ii) and a union bound
make the discarded mass $O_R(M/K)$, uniformly for the actual and stationary
window laws.  For fixed $M,K$ only finitely many words remain; first let
$n\to\infty$, then $K\to\infty$, and finally $M\to\infty$.

The unital self-conjugate algebra generated by finite digit-cylinder
indicators, continuous functions of the real coordinates, and torus
characters separates points on each compact digit truncation of
$\mathcal X_R$.  Stone--Weierstrass therefore makes it uniformly dense there;
the same algebra is convergence determining and its real span is dense in
$L^2(\mu_R)$.  The family $\eta_{n,j}^{(R)}$ is uniformly tight: real and
torus coordinates are compact, and Lemma~\ref{lem:gauss-estimates}(ii)
controls the finitely many digit coordinates.  If \eqref{eq:full-state}
failed, tightness would give a weakly convergent subsequence, say with limit
$\eta$.  Convergence on the determining algebra forces $\eta=\mu_R$, and the
Portmanteau theorem applied to the bounded $\mu_R$-almost-everywhere
continuous function $G$ then contradicts the assumed failure.
\end{proof}

For $z=(x,r,s)$ and $d\in\N_0$, define
\begin{equation}\label{eq:carry-map}
 \phi_z(d)=\lceil x(d+r)-s\rceil.
\end{equation}
Equation \eqref{eq:u-carry} shows that the actual carries satisfy
$d_{j+1}=\phi_{(x_j,\theta_{j-1},\theta_j)}(d_j)$.
If $e=d+r$, the updated real carry obeys
\begin{equation}\label{eq:carry-contraction}
 xe\le e'<xe+1.
\end{equation}
Combining this with \eqref{eq:fibonacci-product}, all actual trajectories and
all pullback trajectories started from zero remain below an absolute integer
$D$.

Set
\begin{equation}\label{eq:reset-set}
 \mathcal R=\left\{(x,r,s):
 x<\frac1{4(D+1)},\quad\frac12<s<\frac34\right\}.
\end{equation}
For $0\le d\le D$ and $z\in\mathcal R$,
\[
 -\frac34<x(d+r)-s<-\frac14,
\]
so $\phi_z(d)=0$.  The set $\mathcal R$ has positive
$\widehat\mu_0$-measure.  By Lemma~\ref{lem:torus-mixing}, almost every
two-sided orbit visits it infinitely often in the past.  Start at time $-N$
with carry zero and iterate to time zero.  After the most recent reset the
answer no longer depends on $N$, so it stabilizes almost surely to a
measurable function $d_*$ satisfying
\begin{equation}\label{eq:carry-graph}
 d_*(\mathcal Sz)=\phi_z(d_*(z)).
\end{equation}
Thus the stationary carry extension is a measurable invariant graph and is
isomorphic to the mixing base--torus system.

For $R\ge1$, define $d_j^{(R)}$ by setting the carry equal to zero at time
$j-R$ and iterating $R$ times.  Let
\[
 p_R=\widehat\mu_0\left(
 \bigcap_{t=0}^{R-1}\mathcal S^{-t}\mathcal R^c\right).
\]
Ergodicity and $\widehat\mu_0(\mathcal R)>0$ imply $p_R\to0$.  Whenever the
actual orbit has a reset during $[j-R,j)$, $d_j=d_j^{(R)}$.  The no-reset
indicator has stationary-null boundary, so Lemma~\ref{lem:full-state} gives
\begin{equation}\label{eq:carry-coupling-prob}
 \sup_{j\in\mathcal J_n}\Pp(d_j\ne d_j^{(R)})\le p_R+o_n(1).
\end{equation}

We now prove the weak law for the remainder in
\eqref{eq:principal-split}.

\begin{proposition}[Bounded-remainder weak law]\label{prop:remainder}
\begin{equation}\label{eq:remainder-WLLN}
 \frac1L\sum_{j\in\mathcal J_n}(-1)^j
 (B_j-\E B_j)\longrightarrow0
 \quad\text{in probability}.
\end{equation}
\end{proposition}

\begin{proof}
Let $B_j^{(R)}$ be the bounded remainder obtained by replacing $d_j$ with
$d_j^{(R)}$.  For fixed $R$, it is a bounded function of a fixed past window,
the torus coordinates, and the current real complete quotient.  Its
discontinuities lie on countably many Gauss boundaries and finitely many
ceiling or fractional-part hypersurfaces, all of stationary measure zero.

By the $L^2$-density established after Lemma~\ref{lem:full-state}, choose
real-valued cylinder--torus polynomials $P_{R,M}$, of finite radii
$R_M\ge R$ that may grow with $M$, such that
\begin{equation}\label{eq:remainder-L2-approx}
 \|B^{(R)}-P_{R,M}\|_{L^2(\mu_{R_M})}
 =:\delta_{R,M}\longrightarrow0.
\end{equation}
They may be taken real-valued because replacing an approximant by its real
part does not increase its distance from the real-valued function $B^{(R)}$.
For each fixed $R,M$, apply Lemma~\ref{lem:full-state} with window radius
$R_M$ to
$|B^{(R)}-P_{R,M}|^2$ to obtain
\begin{equation}\label{eq:actual-L2-transfer}
 \sup_{j\in\mathcal J_n}
 \E|B_j^{(R)}-P_{R,M,j}|^2
 =\delta_{R,M}^2+o_n(1).
\end{equation}
For fixed $R,M$, Proposition~\ref{prop:two-block} shows that the variance of
$\sum_{j\in\mathcal J_n}(-1)^j(P_{R,M,j}-\E P_{R,M,j})$ is $o(L^2)$.
The diagonal contributes $O(L)$; the $O(LH)$ exceptional pairs contribute
$O(LH)=o(L^2)$; and all other pairs contribute $O(L^2\delta_n)=o(1)$, with
$\delta_n$ as in \eqref{eq:delta-n}.  Lemma~\ref{lem:full-state} gives
$\sup_j|\E P_{R,M,j}-\int P_{R,M}\,d\mu_{R_M}|=o(1)$, so replacing the
stationary means in Proposition~\ref{prop:two-block} by the actual means adds
only $o(L^2)$ to the covariance sum.

Write $\widetilde X=X-\E X$.  Minkowski's inequality and
\eqref{eq:actual-L2-transfer} give
\begin{align}\label{eq:Minkowski-polynomial}
 &\left\|\frac1L\sum_{j\in\mathcal J_n}(-1)^j
 (\widetilde B_j^{(R)}-\widetilde P_{R,M,j})\right\|_2\\
 &\qquad\le\frac2L\sum_{j\in\mathcal J_n}
 \|B_j^{(R)}-P_{R,M,j}\|_2
 \le C\delta_{R,M}+o_n(1).\notag
\end{align}
First let $n\to\infty$, then $M\to\infty$.  This proves the analogue of
\eqref{eq:remainder-WLLN} with $B_j^{(R)}$ for every fixed $R$.

Finally, boundedness and \eqref{eq:carry-coupling-prob} imply
\[
 \sup_{j\in\mathcal J_n}\|B_j-B_j^{(R)}\|_2
 \le C\sqrt{p_R+o_n(1)}.
\]
A second centered Minkowski estimate yields
\begin{equation}\label{eq:Minkowski-reset}
 \limsup_{n\to\infty}
 \left\|\frac1L\sum_{j\in\mathcal J_n}(-1)^j
 (\widetilde B_j-\widetilde B_j^{(R)})\right\|_2
 \le C\sqrt{p_R}.
\end{equation}
Letting $R\to\infty$ completes the proof.  The order of limits is
$n\to\infty$, then $M\to\infty$, then $R\to\infty$.
\end{proof}




\section{Stopping, centering, and completion of the proof}

The continuous height has the exact form
\begin{equation}\label{eq:Cq}
 C_j=\frac{n}{q_j+q_{j-1}x_j},
 \qquad N_j=C_j-E_j,\qquad 0\le E_j\le E_*.
\end{equation}

\begin{lemma}[Stopping time and end terms]\label{lem:stopping}
With $H=L^{3/4}$,
\begin{equation}\label{eq:tau}
 \tau_n=\frac L\lambda+O_{\Pp}(H).
\end{equation}
Moreover,
\begin{equation}\label{eq:end-terms}
 \frac1L\left[
 S_n(\alpha)-\sum_{j\in\mathcal J_n}(-1)^j\Phi(x_j,u_j)
 \right]\longrightarrow0
 \quad\text{in probability}.
\end{equation}
\end{lemma}

\begin{proof}
Since $q_j\le q_j+q_{j-1}x_j<2q_j$, \eqref{eq:Cq} gives a deterministic
bracketing.  If $q_j\le n/(2(E_*+1))$, then $C_j\ge E_*+1$ and hence
$N_j\ge1$.  If $q_j>n$, then $C_j<1$ and, because $N_j$ is a nonnegative
integer, $N_j=0$.  Thus $\tau_n$ lies between the hitting times at which
$q_j$ crosses these two fixed multiples of $n$.  Applying
\eqref{eq:q-large-deviation} at $L/\lambda\pm CH$ to both thresholds, and
then choosing $C$ large, proves \eqref{eq:tau}.

Also,
\begin{equation}\label{eq:Phi-tail-bound}
 |\Phi(x,u)|\le\frac1{8x}+\frac12
 \le\frac{a_1(x)}8+\frac58.
\end{equation}
Lemma~\ref{lem:gauss-estimates}(ii) therefore gives
\begin{equation}\label{eq:Phi-tail}
 \sup_{j,n}\Pp(|\Phi(x_j,u_j)|>t)\le\frac{C}{1+t}.
\end{equation}
On the event $|\tau_n-m_n|\le CH$, the indices omitted in
\eqref{eq:end-terms} lie in a deterministic union of $O(H)$ positions.
The probability that one of their costs exceeds $\varepsilon L$ is
$O(H/L)=o(1)$.  The expected sum of the remaining costs, truncated at
$\varepsilon L$, is $O(H\log L)$ by integration of
\eqref{eq:Phi-tail}; divided by $L$ this tends to zero.  Finally let $C$ grow
and use \eqref{eq:tau}.
\end{proof}

\begin{proof}[Proof of Theorem~\ref{thm:main}]
By Proposition~\ref{prop:principal}, Corollary~\ref{cor:principal-cauchy},
and Proposition~\ref{prop:remainder}, there are deterministic $c_n$ such that
\begin{equation}\label{eq:centered-final}
 \frac1L\sum_{j\in\mathcal J_n}(-1)^j\Phi(x_j,u_j)-c_n
 \Longrightarrow \operatorname{Cauchy}\!\left(0,\frac1{2\pi}\right).
\end{equation}
Lemma~\ref{lem:stopping} and Slutsky's theorem replace the bulk sum by
$S_n/L$.

It remains to remove the deterministic centering.  For almost every
$\alpha$,
\[
 \{k(1-\alpha)\}=1-\{k\alpha\},
 \qquad S_n(1-\alpha)=-S_n(\alpha).
\]
Thus $X_n=S_n/\log n$ has an exactly symmetric law.  The family
$X_n-c_n$ is tight by \eqref{eq:centered-final}; symmetry makes
$X_n+c_n$ tight as well.  Since their deterministic center difference is
$2c_n$, the sequence $(c_n)$ is bounded.  If $c_{n_k}\to c$, then
$X_{n_k}\Rightarrow Z+c$, where
$Z\sim\operatorname{Cauchy}(0,1/(2\pi))$.  The limit must be symmetric about
zero, and a nondegenerate translated Cauchy law has this symmetry only when
$c=0$.  Hence every convergent subsequence of $(c_n)$ tends to zero, so
$c_n\to0$.

The Cauchy distribution is continuous; therefore convergence in distribution
is equivalent to the asserted convergence of distribution functions.
\end{proof}

\begin{thebibliography}{99}

\bibitem{Baladi00}
V.~Baladi,
\emph{Positive Transfer Operators and Decay of Correlations},
World Scientific, Singapore, 2000.

\bibitem{Bloom1002}
T.~F. Bloom,
\emph{Erd\H{o}s Problem \#1002},
\url{https://www.erdosproblems.com/1002}, accessed July 13, 2026.

\bibitem{Er64b}
P. Erd\H{o}s,
\emph{Problems and results on diophantine approximations},
Compositio Math. 16 (1964), 52--65.

\bibitem{Ke60}
H. Kesten,
\emph{Uniform distribution mod $1$},
Ann. of Math. (2) 71 (1960), 445--471.

\bibitem{Kallenberg17}
O.~Kallenberg,
\emph{Random Measures, Theory and Applications},
Probability Theory and Stochastic Modelling 77,
Springer, Cham, 2017.

\bibitem{Zygmund59}
A.~Zygmund,
\emph{Trigonometric Series}, 2nd ed.,
Cambridge University Press, Cambridge, 1959.

\end{thebibliography}

\end{document}
